\reportchapter{Modélisation du procédé}
\label{chap:modelisation_procede}

\section{Introduction}
\label{sec:ch3_introduction}

Ce chapitre présente le modèle du procédé utilisé pour valider le BMS en boucle fermée. Il couvre le système batterie, la boucle thermique, les actionneurs de chauffage et de refroidissement ainsi que l'interface véhicule de réponse aux défauts. Les sections suivantes décrivent les données expérimentales, l'identification électrique et thermique, l'implémentation Simulink, puis la validation quantitative du modèle.

\section{Architecture globale de l'environnement de validation}
\label{sec:ch3_architecture}

L'environnement modélisé dans ce chapitre constitue la plateforme complète de simulation du procédé, permettant d'évaluer le BMS en boucle fermée. La vue globale de cette architecture est placée en annexe afin de conserver une lisibilité suffisante sur une page en orientation paysage. Afin de faciliter la lecture, les différents ensembles y sont distingués par un code couleur.

\begin{itemize}
    \item \textbf{Modèle du procédé (en vert) :} ces blocs représentent la dynamique physique des composants externes au logiciel du BMS. Ils comprennent le système batterie, qui regroupe le pack, la boucle de refroidissement, le réchauffeur et le refroidisseur, ainsi que l'interface véhicule de réponse aux défauts.

    \item \textbf{Système sous test (en jaune) :} cet ensemble regroupe l'architecture du système de gestion. Il se compose des modules d'acquisition esclaves, du contrôleur central ainsi que d'une voie de protection matérielle émulant les sécurités physiques indépendantes du logiciel.
\end{itemize}

Les échanges de données entre l'environnement et le BMS reposent sur plusieurs bus de communication, reproduisant l'architecture réseau du véhicule :

\begin{itemize}
    \item \textbf{Bus interne entre maître et esclaves :} il représente conceptuellement le bus CAN interne bidirectionnel. Il permet au maître de requérir des mesures et aux esclaves de transmettre les tensions, températures et diagnostics locaux.

    \item \textbf{Bus véhicule :} il modélise le réseau externe transmettant les informations critiques vers le véhicule. Il véhicule notamment les limitations de puissance calculées par le BMS ainsi que la commande d'ouverture d'urgence des contacteurs.

    \item \textbf{Bus de diagnostic :} il fournit un point d'accès global aux états utiles du BMS maître, constituant une interface essentielle pour l'analyse et la validation de l'algorithmique.
\end{itemize}

Le cœur de la dynamique électrothermique est concentré dans le système batterie. Comme l'illustre la figure~\ref{fig:ch3_battery_system}, ce sous-système reçoit en entrée :
\begin{itemize}
    \item les \textbf{données de mission} en \textbf{vert}, incluant notamment le profil de courant appliqué au pack après correction par l'environnement véhicule ainsi que la température ambiante.
    \item les \textbf{commandes de gestion issues du BMS} en \textbf{jaune}, telles que les commandes thermiques et d'équilibrage.
\end{itemize}

Il fournit en sortie, en \textbf{bleu}, les grandeurs physiques restituées au BMS, en particulier les tensions cellulaires, les températures de zone et les grandeurs globales nécessaires à la supervision.

\begin{figure}[H]
  \centering
  \captionsetup{width=0.75\textwidth, skip=15pt}
  \includegraphics[width=0.9\textwidth]{./figures/vue_globale_systeme_batterie}
  \caption{Vue d'ensemble du système batterie, constituant le cœur physique du modèle de procédé. Les entrées de mission sont représentées en vert, les commandes BMS en jaune et les sorties physiques en bleu.}
  \label{fig:ch3_battery_system}
\end{figure}

\section{Données expérimentales et préparation}
\label{sec:ch3_donnees}

\subsection{Jeu de données exploité}
\label{subsec:ch3_acquisition}

L'identification du modèle s'appuie sur le jeu de données public \emph{Battery and Heating Data in Real Driving Cycles}, disponible sur IEEE DataPort~\cite{Steinstraeter2020_BatteryHeatingData}. Ces données ont été acquises expérimentalement à partir d'un véhicule électrique BMW~i3 instrumenté, circulant dans des conditions réelles de conduite.

\begin{figure}[H]
  \centering
  \captionsetup{width=0.75\textwidth, skip=15pt}
  \includegraphics[width=0.5\textwidth]{./figures/bmw_i3}
  \caption{Véhicule expérimental utilisé pour l'acquisition des données : BMW~i3.}
  \label{fig:ch3_bmw_i3}
\end{figure}

Le jeu de données comprend 70 trajets réels répartis en deux campagnes de mesure complémentaires :

\begin{itemize}
  \item \textbf{Campagne A} : 32 trajets réalisés en été 2019, entre juin et juillet, pour une température ambiante comprise entre 14 et 33{,}5~\textcelsius.
  \item \textbf{Campagne B} : 38 trajets réalisés en hiver 2019--2020, entre décembre et février, pour une température ambiante comprise entre $-3{,}5$ et 14~\textcelsius.
\end{itemize}

Chaque trajet est échantillonné à 10~Hz, soit un pas de temps de 0{,}1~s. Les enregistrements contiennent jusqu'à 48 signaux décrivant le comportement électrique, thermique et dynamique du véhicule : tension et courant du pack, état de charge, température batterie, température ambiante, vitesse, accélération, chauffage et climatisation. Les caractéristiques principales du jeu de données sont synthétisées dans la table~\ref{tab:ch3_dataset}.

\begin{table}[H]
  \centering
  \small
  \begin{tabularx}{\textwidth}{|l|X|X|}
    \hline
    \textbf{Caractéristique} & \textbf{Campagne A} & \textbf{Campagne B} \\
    \hline
    Nombre de trajets & 32 & 38 \\
    \hline
    Nombre de lignes & 467~701 & 627~092 \\
    \hline
    Signaux disponibles & 22--28 & 48 \\
    \hline
    Température ambiante & 14--33{,}5~\textcelsius & $-3{,}5$--14~\textcelsius \\
    \hline
    Température batterie & 16--32~\textcelsius & $-1$--22~\textcelsius \\
    \hline
    Plage de SoC couverte & 34--88{,}5~\% & 0--86{,}1~\% \\
    \hline
    Tension du pack & 349--395~V & 302--394~V \\
    \hline
    Courant du pack & $-395$ à 144~A & $-404$ à 145~A \\
    \hline
    Pas de temps & 0{,}1~s & 0{,}1~s \\
    \hline
  \end{tabularx}
  \captionsetup{width=0.75\textwidth, skip=15pt}
  \caption{Caractéristiques principales des deux campagnes de mesure.}
  \label{tab:ch3_dataset}
\end{table}

La complémentarité des deux campagnes est importante pour l'identification du modèle. La campagne A décrit principalement un fonctionnement thermique nominal. À l'inverse, la campagne B couvre des conditions hivernales, avec des cellules plus froides, des sollicitations plus sévères et une activation plus fréquente du chauffage. Cette diversité permet d'observer le comportement de la batterie sur une large plage de conditions opératoires.

\subsection{Configuration du pack batterie}
\label{subsec:ch3_pack_bmwi3}

Le véhicule étudié est une BMW~i3 de première génération, équipée d'un pack batterie constitué de cellules \emph{Samsung SDI} de 60~Ah. L'architecture du pack est de type \textbf{96s1p} : 96 cellules sont connectées en série et réparties en 8 modules de 12 cellules. Cette configuration conduit à une tension nominale d'environ 355~V et à une énergie utile de l'ordre de 18~kWh~\cite{BatteryDesign_BMWI3, LaPauseGarage_HV2017}.

\begin{figure}[H]
  \centering
  \captionsetup{width=0.75\textwidth, skip=15pt}
  \includegraphics[width=0.5\textwidth]{./figures/battery_pack}
  \caption{Architecture du pack batterie de la BMW~i3.}
  \label{fig:ch3_battery_pack}
\end{figure}

Cette structure 96s1p constitue la base retenue pour la modélisation du pack et pour la définition des interfaces physiques internes.

\subsection{Vérification de la qualité des données}
\label{subsec:ch3_qualite}

Avant l'identification des paramètres, une vérification systématique de la qualité des enregistrements bruts a été réalisée. Cette étape a porté sur la structure des données, la cohérence des en-têtes, la présence de valeurs manquantes, la régularité temporelle de l'échantillonnage et la cohérence physique des signaux. Les principales anomalies détectées sont regroupées dans la table~\ref{tab:ch3_qualite_issues}.

\begin{table}[H]
  \centering
  \small
  \begin{tabularx}{\textwidth}{|l|l|X|}
    \hline
    \textbf{Anomalie} & \textbf{Portée} & \textbf{Description} \\
    \hline
    Espaces dans les noms de colonnes & Toutes les données &
    Présence d'espaces parasites dans certains en-têtes de colonnes. \\
    \hline
    Erreur typographique dans un en-tête & Toutes les données &
    Incohérence typographique dans un en-tête lié à l'état de charge maximal. \\
    \hline
    Colonne dupliquée & Campagne B &
    Une colonne de température liée à la ventilation apparaît deux fois. \\
    \hline
    Valeurs manquantes importantes & Campagne B &
    Un trajet hivernal présente une proportion élevée de valeurs manquantes pour l'état de charge. Deux autres trajets présentent également des lacunes sur le flux de refroidissement. \\
    \hline
    Valeurs manquantes modérées & Campagne A &
    Un trajet estival contient des valeurs absentes pour la température batterie et la température ambiante. \\
    \hline
    Colonnes constantes & Toutes les données &
    Certaines colonnes liées aux bornes de charge et certains indicateurs de chauffage correspondent à des paramètres de configuration
    plutôt qu'à des mesures dynamiques. \\
    \hline
    Défaut d'encodage & Campagne A &
    Présence d'un caractère erroné à la place du symbole de degré, indiquant un problème d'encodage. \\
    \hline
  \end{tabularx}
  \captionsetup{width=0.75\textwidth, skip=15pt}
  \caption{Anomalies identifiées lors de la vérification de la qualité des données.}
  \label{tab:ch3_qualite_issues}
\end{table}

Aucune inversion temporelle, aucun horodatage dupliqué et aucune irrégularité notable d'échantillonnage n'ont été observés. Les principaux signaux physiques, en particulier la tension, le courant et l'état de charge, restent dans des plages cohérentes avec le fonctionnement attendu du véhicule.

\subsection{Nettoyage et consolidation}
\label{subsec:ch3_nettoyage}

Le nettoyage des données a été réalisé par une procédure automatisée afin de produire un jeu consolidé directement exploitable pour l'identification et la validation.

La première étape a consisté à harmoniser les noms de colonnes et à corriger les incohérences relevées lors de l'analyse qualité. Seules les grandeurs communes aux deux campagnes et pertinentes pour la modélisation ont été conservées. Les colonnes constantes ou spécifiques à une seule campagne ont été supprimées afin de garantir une structure homogène.

Un filtrage a ensuite été appliqué pour éliminer les pics isolés correspondant à des variations physiquement irréalistes. Les valeurs manquantes ont été traitées selon la nature des signaux. Les lacunes courtes ont été comblées par interpolation linéaire, tandis que certaines grandeurs, comme les signaux de chauffage ou la vitesse, ont été remplacées par zéro lorsque cette hypothèse restait cohérente avec leur interprétation physique. Les observations contenant encore des valeurs manquantes sur des grandeurs critiques ont finalement été écartées.

Enfin, les grandeurs ont été harmonisées selon une convention commune, et un identifiant de trajet a été ajouté pour distinguer les acquisitions. Après fusion des 70 trajets, le jeu de données nettoyé contient \textbf{1~064~082 lignes} et \textbf{24 colonnes}, sans valeur manquante résiduelle.

\subsection{Analyse exploratoire}
\label{subsec:ch3_eda}

L'analyse exploratoire vise à caractériser les conditions de fonctionnement couvertes par les mesures et à identifier les phénomènes susceptibles d'influencer l'identification des paramètres.

\subsubsection*{Trajet représentatif}

La figure~\ref{fig:ch3_trip_example} présente les principaux signaux d'un trajet hivernal représentatif. Ce trajet comprend une succession de phases dynamiques, une recharge intermédiaire et une reprise de conduite. Pendant la phase de charge, le SoC passe d'environ 33~\% à 71~\%, tandis que la température batterie augmente progressivement d'environ 12~\textcelsius{} à 19~\textcelsius{} dans un environnement proche de 10~\textcelsius.

La figure~\ref{fig:ch3_trip_phases} illustre la segmentation automatique des phases de fonctionnement. Cette segmentation met en évidence les périodes d'accélération, de freinage régénératif, de croisière, de repos, de charge et de chauffage.

\begin{figure}[H]
  \centering
  \captionsetup{width=0.75\textwidth, skip=15pt}
  \includegraphics[width=0.9\textwidth]{./figures/TripB04}
  \caption{Évolution temporelle des principaux signaux d'un trajet hivernal représentatif.}
  \label{fig:ch3_trip_example}
\end{figure}

\begin{figure}[H]
  \centering
  \captionsetup{width=0.75\textwidth, skip=15pt}
  \includegraphics[width=0.9\textwidth]{./figures/TripB04_phases}
  \caption{Segmentation des phases de fonctionnement du trajet hivernal représentatif.}
  \label{fig:ch3_trip_phases}
\end{figure}

\subsubsection*{Écart entre campagnes}

La figure~\ref{fig:ch3_aging_effect} compare les tensions quasi-statiques du pack en fonction du SoC dans la plage thermique commune aux deux campagnes (17--22~\textcelsius). Un décalage systématique apparaît entre les campagnes : la campagne B présente des tensions légèrement plus élevées pour un même SoC.

Cet écart peut traduire une évolution de la relation entre le SoC estimé et la tension à circuit ouvert, liée notamment aux différences de conditions d'utilisation, d'état thermique ou d'état de santé du pack entre les deux campagnes. Afin de limiter l'introduction d'un biais dans l'identification de la courbe OCV et des paramètres électriques, les données des deux campagnes sont donc traitées séparément lors de l'extraction des tables de correspondance.

\begin{figure}[H]
  \centering
  \captionsetup{width=0.75\textwidth, skip=15pt}
  \includegraphics[width=0.9\textwidth]{./figures/aging_effect}
  \caption{Tension quasi-statique du pack en fonction du SoC dans la plage thermique commune aux campagnes A et B.}
  \label{fig:ch3_aging_effect}
\end{figure}

\subsubsection*{Corrélations entre signaux}

La matrice de corrélation de Pearson, présentée à la figure~\ref{fig:ch3_correlation}, met en évidence les relations dominantes entre les principales grandeurs. La tension du pack est fortement corrélée au SoC, tandis que le courant présente une relation plus marquée avec les grandeurs dynamiques du véhicule. La vitesse influence plus indirectement les grandeurs électriques.

Ces observations confirment la pertinence du courant comme entrée principale du modèle électrique équivalent et orientent la sélection des grandeurs utilisées pour l'identification paramétrique.

\begin{figure}[H]
  \centering
  \captionsetup{width=0.75\textwidth, skip=15pt}
  \includegraphics[width=0.8\textwidth]{./figures/correlation}
  \caption{Matrice de corrélation de Pearson des principaux signaux du jeu de données nettoyé.}
  \label{fig:ch3_correlation}
\end{figure}

\section{Identification des paramètres du modèle}
\label{sec:ch3_identification}

\subsection{Identification des tables électriques}
\label{subsec:ch3_luts}

Le comportement électrique de la cellule est décrit par un modèle électrique équivalent à deux branches RC, noté ECM~2RC. Son identification repose sur six tables de correspondance bidimensionnelles, définies au niveau cellule et dépendant du SoC et de la température :
\[
\mathrm{OCV}(\mathrm{SoC},T),\quad
R_0(\mathrm{SoC},T),\quad
R_1(\mathrm{SoC},T),\quad
C_1(\mathrm{SoC},T),\quad
R_2(\mathrm{SoC},T),\quad
C_2(\mathrm{SoC},T).
\]

\subsubsection*{Courbe OCV}

La tension à circuit ouvert, ou \emph{Open-Circuit Voltage} (OCV), est extraite à partir des instants quasi-statiques. Ces instants sont définis par un courant inférieur à 3~A pendant au moins 3~s. Pour chaque bin $(\mathrm{SoC},T)$, la médiane des tensions mesurées est calculée. Les zones insuffisamment couvertes sont ensuite interpolées, puis légèrement lissées afin d'obtenir une surface monotone et physiquement cohérente. Des points d'ancrage issus des bornes usuelles d'une chimie NMC sont également ajoutés pour stabiliser les extrémités de la courbe.

La figure~\ref{fig:ch3_ocv_curves} présente les courbes OCV obtenues ainsi que les médianes extraites des données. Les profils observés sont cohérents avec le comportement attendu d'une cellule lithium-ion de type NMC.

\begin{figure}[H]
  \centering
  \captionsetup{width=0.75\textwidth, skip=15pt}
  \includegraphics[width=0.9\textwidth]{./figures/ocv_curves_with_data}
  \caption{Courbes $\mathrm{OCV}(\mathrm{SoC},T)$ après interpolation et lissage, superposées aux médianes expérimentales et aux points d'ancrage NMC.}
  \label{fig:ch3_ocv_curves}
\end{figure}

\subsubsection*{Paramètres de l'ECM~2RC}

La résistance série $R_0$ est estimée à partir des variations brusques de courant, selon le rapport $|\Delta V|/|\Delta I|$, puis agrégée par médiane dans chaque bin $(\mathrm{SoC},T)$. Les paramètres dynamiques $R_1$, $C_1$, $R_2$ et $C_2$ sont ensuite identifiés par régression non linéaire sur des fenêtres glissantes de 60~s. Les tables OCV et $R_0$ sont alors utilisées comme paramètres fixes. Seules les fenêtres présentant une excitation suffisante et une erreur de reconstruction compatible avec le modèle sont retenues.

Le choix d'un modèle ECM~2RC constitue un compromis entre précision et complexité. Un modèle à trois branches RC peut améliorer localement l'ajustement des données, mais augmente le nombre de paramètres et donc le risque de surajustement. À l'inverse, un modèle trop simple ne représente pas correctement les phénomènes transitoires. Ce compromis entre sous-ajustement, généralisation et surajustement est illustré à la figure~\ref{fig:ch3_bias_variance}.

\begin{figure}[H]
  \centering
  \captionsetup{width=0.75\textwidth, skip=15pt}
  \includegraphics[width=0.8\textwidth]{./figures/fitting}
  \caption{Illustration du compromis entre sous-ajustement, ajustement optimal et surajustement.}
  \label{fig:ch3_bias_variance}
\end{figure}

La figure~\ref{fig:ch3_ecm_luts} présente les tables finales des paramètres électriques. Les tendances obtenues sont physiquement cohérentes : les résistances augmentent aux basses températures et les paramètres dynamiques présentent une variabilité plus marquée selon le SoC.

\begin{figure}[H]
  \centering
  \captionsetup{width=0.75\textwidth, skip=15pt}
  \includegraphics[width=0.9\textwidth]{./figures/ecm_luts}
  \caption{Tables de correspondance des paramètres $R_0$, $R_1$, $C_1$, $R_2$ et $C_2$ en fonction du SoC et de la température.}
  \label{fig:ch3_ecm_luts}
\end{figure}

\subsubsection*{Validation préliminaire des tables électriques}

Une première validation est réalisée avant l'intégration du modèle complet. Pour chaque trajet, les profils mesurés de courant, de SoC et de température sont imposés au modèle ECM~2RC. La tension simulée est ensuite comparée à la tension mesurée.

\begin{table}[H]
  \centering
  \small
  \begin{tabular}{|l|c|c|c|c|}
    \hline
    \textbf{Campagne} & \textbf{Trajets} & \textbf{RMSE médiane} & \textbf{RMSE moyenne} & \textbf{Biais moyen} \\
    \hline
    A -- été & 32 & 7{,}6~mV & 7{,}7~mV & $-3{,}2$~mV \\
    \hline
    B -- hiver & 38 & 13{,}0~mV & 19{,}8~mV & $-9{,}6$~mV \\
    \hline
    \textbf{Global} & \textbf{70} & \textbf{10{,}0~mV} & \textbf{14{,}3~mV} & --- \\
    \hline
  \end{tabular}
  \captionsetup{width=0.75\textwidth, skip=15pt}
  \caption{Validation préliminaire de l'ECM~2RC sur les 70 trajets.}
  \label{tab:ch3_ecm_prevalidation}
\end{table}

Les erreurs restent faibles sur l'ensemble des trajets, avec une RMSE médiane globale de 10~mV. Les écarts plus élevés observés pour la campagne hivernale s'expliquent principalement par une couverture quasi-statique plus limitée aux basses températures. Ces résultats confirment la cohérence des tables électriques et leur pertinence pour l'implémentation du système batterie sous Simulink.

\subsection{Identification des paramètres thermiques}
\label{subsec:ch3_thermique_calibration}

Le modèle thermique repose sur trois paramètres principaux : la capacité thermique effective $C_\text{th}$, le coefficient d'échange avec le liquide de refroidissement $h_\text{cool}$ et le coefficient d'échange avec l'air ambiant $h_\text{amb}$. Ces paramètres sont identifiés par recherche sur grille (\emph{grid search}) en minimisant la RMSE entre la température simulée et la température mesurée. L'évaluation est effectuée sur les trajets d'une durée supérieure à 2\,000~s présentant une excursion thermique d'au moins 2~\textcelsius. Les valeurs retenues sont regroupées dans la table~\ref{tab:ch3_thermal_params}.

\begin{table}[H]
  \centering
  \small
  \begin{tabularx}{\textwidth}{|l|>{\centering\arraybackslash}p{0.28\textwidth}|X|}
    \hline
    \textbf{Paramètre} & \textbf{Valeur identifiée} & \textbf{Interprétation physique} \\
    \hline
    $C_\text{th}$ & 1500~J/K & Capacité thermique effective de la cellule \\
    \hline
    $m_\text{cell}$ & 1{,}5~kg avec $c_p = 1000$~J/(kg$\cdot$K) & Masse thermique équivalente \\
    \hline
    $h_\text{cool}$ & 0{,}030~W/K & Échange thermique avec le liquide de refroidissement \\
    \hline
    $h_\text{amb}$ & 0{,}15~W/K & Échange thermique avec l'air ambiant \\
    \hline
  \end{tabularx}
  \captionsetup{width=0.75\textwidth, skip=15pt}
  \caption{Paramètres thermiques identifiés par recherche sur grille.}
  \label{tab:ch3_thermal_params}
\end{table}

\section{Modélisation électrique du pack}
\label{sec:ch3_electrique}

\subsection{Modèle électrique de la cellule}
\label{subsec:ch3_cellule}

Chaque cellule lithium-ion est représentée par un ECM~2RC. Ce modèle associe une source de tension OCV, une résistance ohmique $R_0$ et deux branches RC permettant de représenter les phénomènes de polarisation transitoire. Le schéma du modèle est présenté à la figure~\ref{fig:ch3_ecm_schema}.

\begin{figure}[H]
  \centering
  \captionsetup{width=0.75\textwidth, skip=15pt}
  \includegraphics[width=0.9\textwidth]{./figures/cell_ecm}
  \caption{Modèle électrique équivalent ECM~2RC d'une cellule lithium-ion.}
  \label{fig:ch3_ecm_schema}
\end{figure}

En adoptant la convention $I>0$ en décharge, la tension terminale de la cellule est donnée par l'équation~\eqref{eq:ch3_vcell} :

\begin{equation}
  V_\text{cell}(t)
  =
  \mathrm{OCV}\!\bigl(\mathrm{SoC},T\bigr)
  -
  I(t)\,R_0\!\bigl(\mathrm{SoC},T\bigr)
  -
  V_{R_1C_1}(t)
  -
  V_{R_2C_2}(t)
  \label{eq:ch3_vcell}
\end{equation}

Les paramètres $\mathrm{OCV}$, $R_0$, $R_1$, $C_1$, $R_2$ et $C_2$ sont évalués à chaque pas de temps par interpolation bilinéaire des tables identifiées à la section~\ref{subsec:ch3_luts}.

\subsubsection*{Dynamique des branches RC}

Pour chaque branche $k \in \{1,2\}$, la tension de polarisation $V_{R_kC_k}$ est décrite par l'équation différentielle~\eqref{eq:ch3_vrc} :

\begin{equation}
  \dot{V}_{R_kC_k}
  =
  -\frac{V_{R_kC_k}}{R_kC_k}
  +
  \frac{I}{C_k}
  \label{eq:ch3_vrc}
\end{equation}

Afin de garantir la stabilité numérique, cette équation est discrétisée à l'aide de sa solution exponentielle exacte, donnée par l'équation~\eqref{eq:ch3_vrc_discrete} :

\begin{equation}
  V_{R_kC_k}[n+1]
  =
  e^{-\Delta t/(R_kC_k)}\,V_{R_kC_k}[n]
  +
  R_k\!\left(1-e^{-\Delta t/(R_kC_k)}\right)I[n]
  \label{eq:ch3_vrc_discrete}
\end{equation}

\begin{figure}[H]
  \centering
  \captionsetup{width=0.75\textwidth, skip=15pt}
  \includegraphics[width=0.9\textwidth]{./figures/rc_dynamics}
  \caption{Implémentation Simulink de la dynamique discrète des branches RC conformément à l'équation~\eqref{eq:ch3_vrc_discrete}.}
  \label{fig:ch3_subsystem_rc}
\end{figure}

\subsubsection*{Propagation de l'état de charge}

L'état de charge est calculé par intégration coulométrique. Pour un SoC exprimé en fraction et une capacité nominale $Q_\text{nom}$ exprimée en ampère-heures, la mise à jour discrète est donnée par l'équation~\eqref{eq:ch3_soc} :

\begin{equation}
  \mathrm{SoC}[n+1]
  =
  \mathrm{SoC}[n]
  -
  \frac{I[n]\Delta t}{3600\,Q_\text{nom}}
  \label{eq:ch3_soc}
\end{equation}

\begin{figure}[H]
  \centering
  \captionsetup{width=0.75\textwidth, skip=15pt}
  \includegraphics[width=0.8\textwidth]{./figures/soc_estimator}
  \caption{Implémentation Simulink de la propagation du SoC par intégration coulométrique selon l'équation~\eqref{eq:ch3_soc}.}
  \label{fig:ch3_subsystem_soc}
\end{figure}

\subsection{Dispersion inter-cellules}
\label{subsec:ch3_dispersion}

Les cellules d'un pack réel ne présentent pas des caractéristiques strictement identiques. Des écarts de capacité, de résistance interne, de température et de SoC initial apparaissent en raison des tolérances de fabrication, des conditions d'utilisation et du vieillissement différentiel.

Pour reproduire ces déséquilibres, des dispersions gaussiennes bornées sont appliquées aux paramètres des cellules. La logique associée à cette dispersion inter-cellules est représentée en vert dans la figure~\ref{fig:ch3_ecm_schema}. Les tirages aléatoires sont générés avec une graine fixe, puis sauvegardés. Cette procédure garantit la reproductibilité des simulations et permet de comparer différentes stratégies de supervision sur une même réalisation du pack.

\subsection{Équilibrage passif des cellules}
\label{subsec:ch3_equilibrage}

L'équilibrage passif piloté par le BMS est représenté par une substitution de courant. Lorsqu'une cellule est sélectionnée pour l'équilibrage, un courant d'équilibrage constant $I_\text{bal}$ est soustrait au courant principal de la cellule selon la convention de signe du modèle. En dehors des phases d'équilibrage, ce courant est nul.

La logique de l'équilibrage passif est mise en évidence en jaune dans la figure~\ref{fig:ch3_cell_balancing}. Cette simplification permet de représenter l'effet énergétique de la résistance de dissipation sans modifier explicitement l'impédance interne de la cellule.

\begin{figure}[H]
  \centering
  \captionsetup{width=0.75\textwidth, skip=15pt}
  \includegraphics[width=0.85\textwidth]{./figures/i_ball}
  \caption{Modélisation du courant d'équilibrage passif appliqué à une cellule.}
  \label{fig:ch3_cell_balancing}
\end{figure}

\subsection{Modélisation du module}
\label{subsec:ch3_module}

Un module est constitué de 12 cellules connectées en série. Sous Simulink, l'implémentation repose sur un mécanisme de réplication du modèle de cellule tout en conservant une interface vectorielle.

Le module reçoit en entrée le vecteur des courants d'équilibrage $I_\text{bal}$ et fournit en sortie les 12 tensions cellulaires ainsi que les SoC associés. L'indexation des vecteurs est fixe. Chaque position correspond à une cellule physique donnée, ce qui permet au BMS de localiser directement les grandeurs mesurées ou estimées.

Sur le plan thermique, la topologie des capteurs est prise en compte par une réduction spatiale. Les 12 températures cellulaires sont regroupées en quatre températures de zone, chacune correspondant à la moyenne de trois cellules adjacentes. Cette représentation est cohérente avec une instrumentation de type un capteur pour trois cellules.

\begin{figure}[H]
  \centering
  \captionsetup{width=0.75\textwidth, skip=15pt}
  \includegraphics[width=0.9\textwidth]{./figures/module}
  \caption{Implémentation Simulink d'un module batterie de 12 cellules en série.}
  \label{fig:ch3_module_simulink}
\end{figure}

\subsection{Modélisation du pack}
\label{subsec:ch3_pack}

Le pack complet est obtenu par l'association en série de 8 modules, soit 96 cellules. Le modèle conserve l'information individuelle de chaque cellule plutôt que de réduire immédiatement les grandeurs à une valeur globale. Les signaux issus des modules sont concaténés selon une indexation fixe à l'échelle du pack.

L'interface transmise au BMS comprend ainsi :
\begin{itemize}
  \item un vecteur de 96 tensions cellulaires.
  \item un vecteur de 96 SoC cellulaires.
  \item un vecteur de 32 températures de zone, correspondant à 8 modules multipliés par 4 zones thermiques.
  \item le courant global du pack.
\end{itemize}

Cette structure vectorielle permet au BMS d'appliquer des fonctions de surveillance, d'estimation et d'équilibrage à l'échelle cellule, module ou pack, sans perte d'information.

\section{Modélisation thermique}
\label{sec:ch3_thermique}

\subsection{Modèle thermique de la cellule}
\label{subsec:ch3_thermique_cellule}

La dynamique thermique de chaque cellule est représentée par un modèle à capacité concentrée. La cellule est assimilée à un nœud thermique unique de température $T_\text{cell}$. Le bilan d'énergie est donné par l'équation~\eqref{eq:ch3_thermique} :

\begin{equation}
  C_\text{th}\,\dot{T}_\text{cell}
  =
  Q_\text{gen}(t)
  -
  h_\text{cool}\,\bigl(T_\text{cell} - T_\text{cool}\bigr)
  -
  h_\text{amb}\,\bigl(T_\text{cell} - T_\text{amb}\bigr)
  \label{eq:ch3_thermique}
\end{equation}

où $C_\text{th}$ est la capacité thermique effective de la cellule, $h_\text{cool}$ le coefficient d'échange avec le fluide caloporteur, $h_\text{amb}$ le coefficient d'échange avec l'air ambiant, $T_\text{cool}$ la température du fluide et $T_\text{amb}$ la température ambiante.

La puissance thermique générée est calculée à partir des pertes ohmiques et des pertes associées aux branches RC, selon l'équation~\eqref{eq:ch3_qgen} :

\begin{equation}
  Q_\text{gen}(t)
  =
  I^2 R_0
  +
  \frac{V_{R_1C_1}^2}{R_1}
  +
  \frac{V_{R_2C_2}^2}{R_2}
  \label{eq:ch3_qgen}
\end{equation}

\begin{figure}[H]
  \centering
  \captionsetup{width=0.75\textwidth, skip=15pt}
  \includegraphics[width=0.85\textwidth]{./figures/thermal_node}
  \caption{Implémentation Simulink du nœud thermique d'une cellule fondée sur les équations~\eqref{eq:ch3_thermique} et~\eqref{eq:ch3_qgen}.}
  \label{fig:ch3_cell_thermal_node}
\end{figure}

\subsection{Boucle de refroidissement}
\label{subsec:ch3_boucle_refroidissement}

La boucle de refroidissement permet de moduler la température du fluide caloporteur appliquée aux cellules. Elle intègre de manière simplifiée deux actionneurs thermiques : le réchauffeur, chargé du réchauffement actif du fluide, et le refroidisseur, assurant son refroidissement par l'intermédiaire du circuit de climatisation.

Dans le modèle de procédé, ces actionneurs sont représentés comme des conditions aux limites dynamiques imposées au fluide. Leurs commandes sont fournies par le BMS, ce qui permet d'évaluer l'effet de la stratégie de gestion thermique sur l'évolution de la température du pack.

\begin{figure}[H]
  \centering
  \captionsetup{width=0.75\textwidth, skip=15pt}
  \includegraphics[width=0.85\textwidth]{./figures/coolant_loop}
  \caption{Schéma Simulink de la boucle de refroidissement du pack.}
  \label{fig:ch3_cooling_loop}
\end{figure}

\section{Modélisation de la réponse du véhicule}
\label{sec:ch3_fault_response}

Dans l'environnement global de validation, le sous-système de réponse véhicule fait partie du modèle du procédé. Il représente de manière simplifiée l'interface entre le système batterie et le reste du véhicule. Son rôle est de répercuter sur le courant effectivement appliqué au pack les limitations de puissance calculées par le BMS ainsi que la commande d'ouverture des contacteurs.

\begin{figure}[H]
  \centering
  \captionsetup{width=0.75\textwidth, skip=15pt}
  \includegraphics[width=0.9\textwidth]{./figures/fault_response}
  \caption{Logique interne simplifiée du sous-système d'interface véhicule.}
  \label{fig:ch3_fault_response}
\end{figure}

Deux fonctions principales sont assurées :
\begin{itemize}
    \item \textbf{Application des limitations de puissance :} à partir de la tension du pack et des puissances limites calculées par le BMS en charge et en décharge, le bloc déduit les bornes de courant admissibles. Le courant demandé par le véhicule est alors corrigé afin de respecter ces contraintes avant d'être appliqué au système batterie.

    \item \textbf{Ouverture des contacteurs :} lorsqu'une commande de mise en sécurité est émise par le BMS, le bloc force l'ouverture virtuelle des contacteurs principaux. Le courant appliqué au pack devient alors nul, ce qui permet de reproduire l'isolement électrique de la batterie dans la simulation.
\end{itemize}

\section{Validation du modèle de procédé}
\label{sec:ch3_validation}

\subsection{Méthodologie de validation}
\label{subsec:ch3_strategie_validation}

La validation présentée ici se concentre sur le cœur électrothermique de l'environnement, à savoir le système batterie, afin d'évaluer sa capacité à reproduire fidèlement les mesures expérimentales acquises lors des 70 trajets disponibles. Le bloc de réponse véhicule est conservé dans l'environnement global, mais sa logique dépend surtout des commandes BMS et sera donc exploitée dans les validations du système sous test. La validation du modèle physique est conduite en deux étapes complémentaires.

La première étape correspond à une validation électrique isolée. Le modèle ECM~2RC est alimenté par les profils mesurés de courant, de SoC et de température. Cette configuration permet d'évaluer directement la qualité des tables électriques, indépendamment des erreurs de propagation du SoC et de la température.

La seconde étape correspond à une validation globale en boucle fermée. Le modèle complet, électrique et thermique, est simulé de manière autonome à partir du courant et des conditions initiales. Le SoC et la température utilisés par le modèle sont alors ceux propagés par la simulation. La dispersion inter-cellules est désactivée dans cette étape afin d'évaluer les performances du modèle nominal.

\subsection{Résultats quantitatifs}
\label{subsec:ch3_resultats_convergence}

Les performances globales sont évaluées à partir de trois grandeurs : la tension pack $V_\text{pack}$, le SoC pack $\mathrm{SoC}_\text{pack}$ et la température pack $T_\text{pack}$. Les résultats obtenus sur les 70 trajets sont présentés dans la table~\ref{tab:ch3_convergence_results}.

\begin{table}[H]
  \centering
  \small
  \begin{tabular}{|l|c|c|c|c|}
    \hline
    \textbf{Grandeur} & \textbf{RMSE médiane} & \textbf{RMSE moyenne} & \textbf{RMSE max} & \textbf{Seuil visé} \\
    \hline
    $V_\text{pack}$ & 1{,}045~V & 1{,}482~V & 6{,}327~V & 4{,}0~V \\
    \hline
    $\mathrm{SoC}_\text{pack}$ & 0{,}445~\% & 0{,}498~\% & 1{,}829~\% & 1{,}0~\% \\
    \hline
    $T_\text{pack}$ & 0{,}702~\textcelsius & 0{,}904~\textcelsius & 3{,}570~\textcelsius & 1{,}0~\textcelsius \\
    \hline
  \end{tabular}
  \captionsetup{width=0.75\textwidth, skip=15pt}
  \caption{Métriques de validation du modèle de procédé sur les 70 trajets.}
  \label{tab:ch3_convergence_results}
\end{table}

Les RMSE médianes respectent les seuils visés pour les trois grandeurs. La tension pack présente une marge importante par rapport au seuil de 4~V, tandis que l'erreur médiane de SoC reste inférieure à 0{,}5~\%. La température pack présente également une erreur médiane inférieure au seuil de 1~\textcelsius.

La figure~\ref{fig:ch3_convergence_tripA02} illustre la convergence du modèle sur un trajet représentatif. Les grandeurs simulées suivent correctement les mesures expérimentales, y compris lors des principales phases transitoires.

\begin{figure}[H]
  \centering
  \captionsetup{width=0.75\textwidth, skip=15pt}
  \includegraphics[width=0.9\textwidth]{./figures/TripA02}
  \caption{Comparaison des grandeurs simulées et mesurées sur un trajet représentatif.}
  \label{fig:ch3_convergence_tripA02}
\end{figure}

Les erreurs maximales observées correspondent à des écarts ponctuels lors de transitoires sévères ou dans des zones moins couvertes par les données d'identification. La marge plus réduite sur la température s'explique par les limites du modèle thermique à capacité concentrée, qui ne représente pas les gradients spatiaux internes au module. Ce compromis demeure néanmoins compatible avec le niveau de précision recherché et avec les contraintes de coût de calcul associées aux simulations embarquées.

\section{Conclusion}
\label{sec:ch3_conclusion}

Ce chapitre a présenté la modélisation du procédé utilisé pour l'évaluation du BMS. L'identification du modèle électrique ECM~2RC et le calibrage du modèle thermique ont permis d'obtenir une représentation fiable du système batterie sur une large plage de fonctionnement. L'architecture mise en place replace ce cœur physique dans un environnement cohérent avec les interactions véhicule et la réponse aux défauts.

Les résultats de validation confirment que les erreurs sur la tension, l'état de charge et la température demeurent compatibles avec les objectifs de précision fixés. L'environnement complet ainsi établi constitue une plateforme robuste pour le développement, l'analyse et l'évaluation des algorithmes du système de gestion de batterie, dont les logiques de contrôle et de surveillance seront présentées dans le chapitre suivant.